Assume first that $\operatorname{char}k\ne2$. Eliminating $w$ gives $y^2z-x^3+xz^2=0$. On $y\ne0$, the inverse of projection is $(x:y:z)\mapsto(x:y:z:(x^2-xz)/y)$; on $x+z\ne0$, it is $(x:y:z)\mapsto(x:y:z:yz/(x+z))$. They agree on the overlap by the cubic equation. The only point of the cubic outside these two opens is $(1:0:-1)$. A point of $Y$ with $y=x+z=0$ satisfies $2x^2=0$, so it is $P$. This proves the claimed isomorphism.

On $z\ne0$ the cubic is $y^2=x^3-x$, whose singularity equations would require $y=0$, $3x^2=1$, and $x(x^2-1)=0$, which are incompatible. Its only point on $z=0$ is $(0:1:0)$, where the derivative with respect to $z$ is $1$. Thus the cubic is nonsingular and irreducible by \exref{I.5.9}.

In the chart $w=1$, the two equations of $Y$ have independent linear terms $-y$ and $-x-z$ at $P$. Hence $Y$ is nonsingular of dimension $1$ at $P$. Since $Y-P$ is an irreducible curve and $P$ is not an isolated point, its closure is all of $Y$. Thus $Y$ is irreducible and nonsingular.

In characteristic $2$ the assertion fails: $Y$ contains the line $y=0$, $x=z$, whose points other than $P$ all project to $(1:0:1)$. Moreover, at $(1:0:1:1)$ the two Jacobian rows coincide, so $Y$ is singular there.
